\documentclass[10pt]{extarticle}
\usepackage[a4paper,margin=1.5cm]{geometry}
\usepackage[T1]{fontenc}
\usepackage{amsmath,amssymb}
\usepackage{graphicx}
\usepackage[export]{adjustbox}
\usepackage{multicol}
\usepackage{enumitem}
\usepackage{booktabs,array,tabularx}
\setlength{\columnsep}{0.7cm}
\setlength{\parindent}{0pt}
\setlength{\parskip}{3pt}
\pagestyle{empty}
\makeatletter
% figure registry (figdefs.tex is generated by the extract step)
\newcommand{\qfigdef}[4]{\@namedef{qf@file@#1}{#2}\@namedef{qf@w@#1}{#3}\@namedef{qf@wide@#1}{#4}}
\newcommand{\qf@check}[1]{\@ifundefined{qf@file@#1}{\PackageError{score2}{Unknown figure id #1}{}}{}}
\newcommand{\qf@img}[4][]{\includegraphics[width=#2\dimexpr\@nameuse{qf@w@#4}pt\relax,max width=#3,max height=0.45\textheight,#1]{\@nameuse{qf@file@#4}}}
% Full-width questions: finalize wraps every question that contains a wide figure in \qfullbegin...\qfullend,
% so the WHOLE question (stem, figures, options) is set across the page in reading order, then the two columns
% resume. The question's head - its number, stem and first figure (\qfullsplit marks the end of the head) -
% is kept on one page, so a number is never left at a page foot with its figure overleaf; the rest may break
% across pages normally, so no big gaps are left. Without a \qfullsplit, the whole question is kept together
% if it is under 45% of a page. Never write these yourself.
\newif\ifqf@full
\newif\ifqf@boxed
\newsavebox{\qf@box}
\newcommand{\qf@open}{\global\qf@fulltrue\par\global\qf@boxedtrue
  \setbox\qf@box\vbox\bgroup\hsize\textwidth\linewidth\textwidth}
\newcommand{\qf@emit}[1]{\par\egroup\global\qf@boxedfalse
  \ifdim\dimexpr\ht\qf@box+\dp\qf@box\relax<#1\textheight\noindent\usebox{\qf@box}\par\else\unvbox\qf@box\fi}
\newcommand{\qfullsplit}{\ifqf@boxed\qf@emit{0.85}\fi}
\newcommand{\qf@close}{\ifqf@boxed\qf@emit{0.45}\fi\global\qf@fullfalse}
\newcommand{\qfullbegin}{\end{multicols}\qf@open}
\newcommand{\qfullend}{\qf@close\begin{multicols}{2}}
\newcommand{\qfullstart}{\qf@open}
\newcommand{\qfulldone}{\qf@close}
% \qfig[scale]{id}: block figure, centred, directly where it is written (size set by the tool).
% Fallback if a wide figure is not inside a full-width question: break out of the columns in place - no floats.
\newcommand{\qfig}[2][1]{\qf@check{#2}%
  \ifqf@full\qf@blk{#1}{#2}\else
  \ifnum\@nameuse{qf@wide@#2}=1 \expandafter\@firstoftwo\else\expandafter\@secondoftwo\fi
  {\end{multicols}\begin{center}\qf@img{#1}{\textwidth}{#2}\end{center}\begin{multicols}{2}}%
  {\qf@blk{#1}{#2}}\fi}
\newcommand{\qf@blk}[2]{\par\smallskip{\centering\qf@img{#1}{\linewidth}{#2}\par}\smallskip}
% \qfigi[scale]{id}: inline figure (inside option lists / tables / minipages); never breaks columns.
\newcommand{\qfigi}[2][1]{\qf@check{#2}\qf@img[valign=c]{#1}{\linewidth}{#2}}
\makeatother
% structure macros - lint reads these to check completeness and figure placement
\newcommand{\papersec}[1]{\par\bigskip{\centering\Large\bfseries PAPER-#1\par}\medskip}
\newcommand{\partsec}[2]{\par\medskip\noindent{\large\bfseries PART-#1: #2}\par\smallskip}
\newcommand{\sect}[1]{\par\smallskip\noindent\textbf{#1}\par}
\newcommand{\qn}[1]{\par\smallskip\noindent\textbf{#1.}~}
\newcommand{\paraq}[2]{\par\smallskip\noindent\textbf{Paragraph for Questions #1 and #2}\par}
\newenvironment{qoptions}{\begin{itemize}[leftmargin=1.8em,itemsep=1pt,topsep=2pt,label={}]}{\end{itemize}}
\newcommand{\opt}[2]{\item[(#1)] #2}
% figure registry (generated by score2_tools.py finalize)
\qfigdef{p1_phy_s1_q03_1}{Score Advanced - Test 18 (27 May 2023)-img/p1_phy_s1_q03_1.png}{153.3}{0}
\qfigdef{p1_phy_s1_q04_1}{Score Advanced - Test 18 (27 May 2023)-img/p1_phy_s1_q04_1.png}{128.9}{0}
\qfigdef{p1_phy_s2_q07_1}{Score Advanced - Test 18 (27 May 2023)-img/p1_phy_s2_q07_1.png}{129.6}{0}
\qfigdef{p1_phy_s2_q08_1}{Score Advanced - Test 18 (27 May 2023)-img/p1_phy_s2_q08_1.png}{140.4}{0}
\qfigdef{p1_phy_s2_q09_1}{Score Advanced - Test 18 (27 May 2023)-img/p1_phy_s2_q09_1.png}{113.6}{0}
\qfigdef{p1_phy_s2_q09_2}{Score Advanced - Test 18 (27 May 2023)-img/p1_phy_s2_q09_2.png}{102.6}{0}
\qfigdef{p1_phy_s2_q09_3}{Score Advanced - Test 18 (27 May 2023)-img/p1_phy_s2_q09_3.png}{103.1}{0}
\qfigdef{p1_phy_s2_q09_4}{Score Advanced - Test 18 (27 May 2023)-img/p1_phy_s2_q09_4.png}{103.1}{0}
\qfigdef{p1_phy_s2_q09_5}{Score Advanced - Test 18 (27 May 2023)-img/p1_phy_s2_q09_5.png}{103.1}{0}
\qfigdef{p1_phy_s2_q10_1}{Score Advanced - Test 18 (27 May 2023)-img/p1_phy_s2_q10_1.png}{69.6}{0}
\qfigdef{p1_phy_s3_q06_1}{Score Advanced - Test 18 (27 May 2023)-img/p1_phy_s3_q06_1.png}{83.2}{0}
\qfigdef{p1_phy_s3_q07_1}{Score Advanced - Test 18 (27 May 2023)-img/p1_phy_s3_q07_1.png}{97.9}{0}
\qfigdef{p1_chem_s1_q02_1}{Score Advanced - Test 18 (27 May 2023)-img/p1_chem_s1_q02_1.png}{192.0}{0}
\qfigdef{p1_chem_s1_q04_A}{Score Advanced - Test 18 (27 May 2023)-img/p1_chem_s1_q04_A.png}{147.0}{0}
\qfigdef{p1_chem_s1_q04_B}{Score Advanced - Test 18 (27 May 2023)-img/p1_chem_s1_q04_B.png}{166.6}{0}
\qfigdef{p1_chem_s1_q04_C}{Score Advanced - Test 18 (27 May 2023)-img/p1_chem_s1_q04_C.png}{187.4}{0}
\qfigdef{p1_chem_s1_q04_D}{Score Advanced - Test 18 (27 May 2023)-img/p1_chem_s1_q04_D.png}{149.9}{0}
\qfigdef{p1_chem_s1_q05_1}{Score Advanced - Test 18 (27 May 2023)-img/p1_chem_s1_q05_1.png}{192.8}{0}
\qfigdef{p1_chem_s1_q06_A}{Score Advanced - Test 18 (27 May 2023)-img/p1_chem_s1_q06_A.png}{74.7}{0}
\qfigdef{p1_chem_s1_q06_B}{Score Advanced - Test 18 (27 May 2023)-img/p1_chem_s1_q06_B.png}{104.4}{0}
\qfigdef{p1_chem_s1_q06_C}{Score Advanced - Test 18 (27 May 2023)-img/p1_chem_s1_q06_C.png}{74.5}{0}
\qfigdef{p1_chem_s1_q06_D}{Score Advanced - Test 18 (27 May 2023)-img/p1_chem_s1_q06_D.png}{84.6}{0}
\qfigdef{p1_chem_s2_q09_1}{Score Advanced - Test 18 (27 May 2023)-img/p1_chem_s2_q09_1.png}{53.6}{0}
\qfigdef{p2_phy_s1_q01_1}{Score Advanced - Test 18 (27 May 2023)-img/p2_phy_s1_q01_1.png}{146.9}{0}
\qfigdef{p2_phy_s1_q02_1}{Score Advanced - Test 18 (27 May 2023)-img/p2_phy_s1_q02_1.png}{82.0}{0}
\qfigdef{p2_phy_s1_q03_1}{Score Advanced - Test 18 (27 May 2023)-img/p2_phy_s1_q03_1.png}{107.8}{0}
\qfigdef{p2_phy_s1_q04_1}{Score Advanced - Test 18 (27 May 2023)-img/p2_phy_s1_q04_1.png}{108.3}{0}
\qfigdef{p2_phy_s1_q04_A}{Score Advanced - Test 18 (27 May 2023)-img/p2_phy_s1_q04_A.png}{65.4}{0}
\qfigdef{p2_phy_s1_q04_B}{Score Advanced - Test 18 (27 May 2023)-img/p2_phy_s1_q04_B.png}{76.5}{0}
\qfigdef{p2_phy_s1_q04_C}{Score Advanced - Test 18 (27 May 2023)-img/p2_phy_s1_q04_C.png}{73.2}{0}
\qfigdef{p2_phy_s1_q04_D}{Score Advanced - Test 18 (27 May 2023)-img/p2_phy_s1_q04_D.png}{84.4}{0}
\qfigdef{p2_phy_s2_q05_1}{Score Advanced - Test 18 (27 May 2023)-img/p2_phy_s2_q05_1.png}{205.4}{0}
\qfigdef{p2_phy_s2_q06_1}{Score Advanced - Test 18 (27 May 2023)-img/p2_phy_s2_q06_1.png}{122.0}{0}
\qfigdef{p2_phy_s2_q07_1}{Score Advanced - Test 18 (27 May 2023)-img/p2_phy_s2_q07_1.png}{117.5}{0}
\qfigdef{p2_phy_s2_q08_1}{Score Advanced - Test 18 (27 May 2023)-img/p2_phy_s2_q08_1.png}{207.6}{0}
\qfigdef{p2_phy_s2_q09_1}{Score Advanced - Test 18 (27 May 2023)-img/p2_phy_s2_q09_1.png}{205.8}{0}
\qfigdef{p2_phy_s2_q10_1}{Score Advanced - Test 18 (27 May 2023)-img/p2_phy_s2_q10_1.png}{30.4}{0}
\qfigdef{p2_phy_s3_q01_1}{Score Advanced - Test 18 (27 May 2023)-img/p2_phy_s3_q01_1.png}{108.2}{0}
\qfigdef{p2_phy_s3_q06_1}{Score Advanced - Test 18 (27 May 2023)-img/p2_phy_s3_q06_1.png}{91.4}{0}
\qfigdef{p2_chem_s1_q02_B}{Score Advanced - Test 18 (27 May 2023)-img/p2_chem_s1_q02_B.png}{65.9}{0}
\qfigdef{p2_chem_s1_q02_C}{Score Advanced - Test 18 (27 May 2023)-img/p2_chem_s1_q02_C.png}{77.7}{0}
\qfigdef{p2_chem_s1_q02_D}{Score Advanced - Test 18 (27 May 2023)-img/p2_chem_s1_q02_D.png}{89.2}{0}
\qfigdef{p2_chem_s1_q03_1}{Score Advanced - Test 18 (27 May 2023)-img/p2_chem_s1_q03_1.png}{99.4}{0}
\qfigdef{p2_chem_s1_q03_A}{Score Advanced - Test 18 (27 May 2023)-img/p2_chem_s1_q03_A.png}{47.0}{0}
\qfigdef{p2_chem_s1_q03_B}{Score Advanced - Test 18 (27 May 2023)-img/p2_chem_s1_q03_B.png}{39.6}{0}
\qfigdef{p2_chem_s1_q03_C}{Score Advanced - Test 18 (27 May 2023)-img/p2_chem_s1_q03_C.png}{46.6}{0}
\qfigdef{p2_chem_s1_q03_D}{Score Advanced - Test 18 (27 May 2023)-img/p2_chem_s1_q03_D.png}{41.1}{0}
\qfigdef{p2_chem_s2_q08_1}{Score Advanced - Test 18 (27 May 2023)-img/p2_chem_s2_q08_1.png}{181.8}{0}
\qfigdef{p2_chem_s2_q08_A}{Score Advanced - Test 18 (27 May 2023)-img/p2_chem_s2_q08_A.png}{44.8}{0}
\qfigdef{p2_chem_s2_q08_B}{Score Advanced - Test 18 (27 May 2023)-img/p2_chem_s2_q08_B.png}{27.0}{0}
\qfigdef{p2_chem_s2_q08_C}{Score Advanced - Test 18 (27 May 2023)-img/p2_chem_s2_q08_C.png}{52.5}{0}
\qfigdef{p2_chem_s2_q08_D}{Score Advanced - Test 18 (27 May 2023)-img/p2_chem_s2_q08_D.png}{96.7}{0}
\qfigdef{p2_chem_s2_q09_1}{Score Advanced - Test 18 (27 May 2023)-img/p2_chem_s2_q09_1.png}{289.9}{1}
\qfigdef{p2_chem_s2_q09_A}{Score Advanced - Test 18 (27 May 2023)-img/p2_chem_s2_q09_A.png}{71.5}{0}
\qfigdef{p2_chem_s2_q09_B}{Score Advanced - Test 18 (27 May 2023)-img/p2_chem_s2_q09_B.png}{41.1}{0}
\qfigdef{p2_chem_s2_q09_C}{Score Advanced - Test 18 (27 May 2023)-img/p2_chem_s2_q09_C.png}{71.3}{0}
\qfigdef{p2_chem_s2_q09_D}{Score Advanced - Test 18 (27 May 2023)-img/p2_chem_s2_q09_D.png}{75.6}{0}
\qfigdef{p2_chem_s3_q05_1}{Score Advanced - Test 18 (27 May 2023)-img/p2_chem_s3_q05_1.png}{236.3}{0}
\qfigdef{p2_chem_s3_q05_2}{Score Advanced - Test 18 (27 May 2023)-img/p2_chem_s3_q05_2.png}{203.3}{0}
\qfigdef{p2_chem_s3_q06_1}{Score Advanced - Test 18 (27 May 2023)-img/p2_chem_s3_q06_1.png}{190.6}{0}
\qfigdef{p2_chem_s3_q08_1}{Score Advanced - Test 18 (27 May 2023)-img/p2_chem_s3_q08_1.png}{238.8}{0}
\qfigdef{p2_math_s1_q04_1}{Score Advanced - Test 18 (27 May 2023)-img/p2_math_s1_q04_1.png}{59.3}{0}
\raggedcolumns
\makeatletter\@beginparpenalty=10000\makeatother
\newcolumntype{L}[1]{>{\raggedright\arraybackslash}p{#1}}
\begin{document}
\begin{center}{\Large\bfseries Score Advanced --- Test 18 (27 May 2023)}\end{center}
\begin{multicols}{2}
\raggedcolumns
\papersec{1}
\partsec{1}{PHYSICS}
\sect{SECTION-I (i)}
\qn{1} The velocity, acceleration and force in two systems of units are related to each other as under\par
(i) $v'=\frac{\alpha^2}{\beta}v$\quad (ii) $a'=(\alpha\beta)a$\quad (iii) $F'=\left(\frac{1}{\alpha\beta}\right)F$\par
where, all the primed symbols belong to one system of units and the unprimed symbols belong to the other system of units. Here $\alpha$ and $\beta$ are dimensionless constants. Which of the following is / are correct ?
\begin{qoptions}
\opt{A}{The standard of length in each of these two systems are related to each other as $\ell'=\frac{\alpha^3}{\beta^3}\ell$.}
\opt{B}{The standard of mass in each of these two systems are related to each other as $m'=\left(\frac{1}{\alpha^2\beta^2}\right)m$.}
\opt{C}{The standard of time in each of these two systems are related to each other as $t'=\left(\frac{\alpha}{\beta^2}\right)t$.}
\opt{D}{The standard of momentum in each of these two systems are related to each other as $p'=\left(\frac{1}{\beta^3}\right)p$.}
\end{qoptions}
\qn{2} A cubical box of side 1 meter contains helium gas (atomic weight 4) at a pressure of 100 Nm$^{-2}$. During an observation time of 1 second, it was found there are 500 collisions per helium atom perpendicular to wall. Ideal gas theory is valid in this given situation. Take $R=\frac{25}{3}$ J mol$^{-1}$k$^{-1}$ and $k=1.38\times10^{-23}$ JK$^{-1}$. Choose the correct options.
\begin{qoptions}
\opt{A}{Temperature of the gas is 160 K.}
\opt{B}{Average K.E. per atom is $3.312\times10^{-21}$ J.}
\opt{C}{Total mass of helium gas in the box is $3\times10^{-4}$ kg.}
\opt{D}{Total mass of helium gas in the box is $2\times10^{-4}$ kg.}
\end{qoptions}
\qn{3} A composite block is made of slabs A, B, C, D and E of different thermal conductivities (given in terms of a constant K) and sizes (given in terms of length L) as shown in the figure. All slabs are of same width. Heat $Q$ flows only from left to right through the blocks. Then in steady state :-\par
\qfig{p1_phy_s1_q03_1}
\begin{qoptions}
\opt{A}{heat flow rate through A and E slabs are same.}
\opt{B}{heat flow rate through slab C is more than rate of heat flow through B or D.}
\opt{C}{temperature difference across slab E is smallest.}
\opt{D}{Heat flow rate through C = Heat flow rate through B + Heat flow rate through D}
\end{qoptions}
\qn{4} A particle A of mass $m$ moving on a smooth horizontal surface collides head on with a stationary particle B of mass $2m$ which is situated at a distance $d$ from a wall. The coefficient of restitution between A \& B \& between B \& wall is $e=1/4$. (Neglect size of particles in comparison to $d$)\par
\qfig{p1_phy_s1_q04_1}
\begin{qoptions}
\opt{A}{Velocity of ball A just after first collision is $v/6$.}
\opt{B}{Velocity of ball B just after first collision is $5v/12$.}
\opt{C}{Distance from the wall where they collide again if entire motion takes place along a straight line perpendicular to wall is $\frac{3d}{13}$.}
\opt{D}{Distance from the wall where they collide again if entire motion takes place along a straight line perpendicular to wall is $\frac{2d}{13}$.}
\end{qoptions}
\qn{5} In electromagnetic radiation, the variation of electric field with time at a particular position is given by\par
$E=E_m\left[2\cos^2\left(5\pi\times10^{15}t\right)-1\right]$ where $t$ is in second.\par
Above radiation is absorbed by a sample of gas having atomic hydrogen in ground state. Which of the following options are correct ? [$h=6.6\times10^{-34}$ J$-$s]
\begin{qoptions}
\opt{A}{The energy of photon in the radiation is 10.3 eV}
\opt{B}{Ionization of hydrogen atom is not possible by this radiation.}
\opt{C}{The approximate value of ratio of kinetic energy of e$^{-}$ in hydrogen atom in ground state to kinetic energy of e$^{-}$ in ionized state is 1.9 : 1}
\opt{D}{Stopping potential of 7V (approximately) is required to stop the electron knocked out from hydrogen atom by above radiation.}
\end{qoptions}
\qn{6} A negative charged particle with velocity $\vec v=y\hat j+z\hat k$ moves in a magnetic field $\vec B=z\hat j+y\hat k$. The magnitude of magnetic force on the particle is $F$. Choose the wrong options.
\begin{qoptions}
\opt{A}{$F=0$ if $y=2z$.}
\opt{B}{$F$ is always along +ve $x$ axis for any value of $y$ and $z$.}
\opt{C}{$F$ is along +ve $x$ axis if $z>y$.}
\opt{D}{$F$ is along $-$ve $x$ axis if $z=y$.}
\end{qoptions}
\sect{SECTION-I (ii)}
\qn{7} The circuit below has an ideal inductor and a switch as shown in figure. List-I gives the different circuit elements and list-II gives the value.\par
\qfig{p1_phy_s2_q07_1}
{\footnotesize\setlength{\tabcolsep}{3pt}\noindent\begin{tabular}{|c|L{3.3cm}|c|c|}\hline
\multicolumn{2}{|c|}{\textbf{List-I}} & \multicolumn{2}{c|}{\textbf{List-II}}\\\hline
(I) & Current in inductor & (P) & $\dfrac{6V}{7}$\\\hline
(II) & Voltage across inductor & (Q) & $\dfrac{3V}{2}$\\\hline
(III) & Current supplied by battery & (R) & $\dfrac{2V}{7R}$\\\hline
(IV) & Energy stored in the inductor & (S) & $\dfrac{V}{R}$\\\hline
 & & (T) & Zero\\\hline
\end{tabular}}\par
If the switch was closed at $t=0$ then which of the following options are correct match of list I and II for that instant ?
\begin{qoptions}
\opt{A}{I $\rightarrow$ T;II $\rightarrow$ P;III $\rightarrow$ R;IV $\rightarrow$ T}
\opt{B}{I $\rightarrow$ P;II $\rightarrow$ Q;III $\rightarrow$ R;IV $\rightarrow$ S}
\opt{C}{I $\rightarrow$ T;II $\rightarrow$ R;III $\rightarrow$ P;IV $\rightarrow$ T}
\opt{D}{I $\rightarrow$ R;II $\rightarrow$ S;III $\rightarrow$ T;IV $\rightarrow$ P}
\end{qoptions}
\qn{8} Four combination of two thin lenses/lens and mirror are given in list-I. The radius of curvature of all curved surfaces is same and the refractive index of all the lenses is 1.5. Focal length of a equiconvex lens kept in air with same radius and same refractive index as that of surface here is $f$. Match lens combination in list-I with their modulus of equivalent focal length in air in list-II and select the correct answer using the code given below the lists.\par
\qfig{p1_phy_s2_q08_1}
\begin{qoptions}
\opt{A}{I $\rightarrow$ R;II $\rightarrow$ Q;III $\rightarrow$ S;IV $\rightarrow$ P}
\opt{B}{I $\rightarrow$ R;II $\rightarrow$ Q;III $\rightarrow$ R;IV $\rightarrow$ P}
\opt{C}{I $\rightarrow$ S;II $\rightarrow$ Q;III $\rightarrow$ R;IV $\rightarrow$ S}
\opt{D}{I $\rightarrow$ Q;II $\rightarrow$ S;III $\rightarrow$ P;IV $\rightarrow$ R}
\end{qoptions}
\qn{9} A train is moving with speed $u$. Engine of train blows a long whistle. The sound of whistle is composed of components varying from $f_1$ to $f_2$ as shown in figure. List-I shows the distribution of the sound intensity of the whistle as observed by different observers. Match List-I with List-II.\par
\qfig{p1_phy_s2_q09_1}
{\footnotesize\setlength{\tabcolsep}{2pt}\noindent\begin{tabular}{|c|c|c|L{2.9cm}|}\hline
\multicolumn{2}{|c|}{\textbf{List-I}} & \multicolumn{2}{c|}{\textbf{List-II}}\\\hline
(I) & \rule{0pt}{2.2cm}\qfigi[0.75]{p1_phy_s2_q09_2} & (P) & Observer in train\\\hline
(II) & \rule{0pt}{2.2cm}\qfigi[0.75]{p1_phy_s2_q09_3} & (Q) & Observer on ground close to track and infront of train\\\hline
(III) & \rule{0pt}{2.2cm}\qfigi[0.75]{p1_phy_s2_q09_4} & (R) & Observer on ground close to track and behind train\\\hline
(IV) & \rule{0pt}{2.2cm}\qfigi[0.75]{p1_phy_s2_q09_5} & (S) & Observer on ground for which engine approaches and then recedes\\\hline
 & & (T) & None of these\\\hline
\end{tabular}}\par
\begin{qoptions}
\opt{A}{I $\rightarrow$ Q;II $\rightarrow$ R;III $\rightarrow$ S;IV $\rightarrow$ P}
\opt{B}{I $\rightarrow$ Q;II $\rightarrow$ R;III $\rightarrow$ P;IV $\rightarrow$ S}
\opt{C}{I $\rightarrow$ S;II $\rightarrow$ Q;III $\rightarrow$ P;IV $\rightarrow$ R}
\opt{D}{I $\rightarrow$ R;II $\rightarrow$ P;III $\rightarrow$ S;IV $\rightarrow$ Q}
\end{qoptions}
\qn{10} Answer the following by appropriately matching the lists based on the information given in the paragraph.\par
A block P of mass $m$ is placed on smooth horizontal surface and connected with a spring. Another block Q of same mass is placed on the block P. The two blocks are pulled by a distance A and released. The blocks oscillates without slipping on each other. Then match the matrix.\par
\qfig{p1_phy_s2_q10_1}
{\footnotesize\setlength{\tabcolsep}{2pt}\noindent\begin{tabular}{|c|L{1.6cm}|c|L{3.7cm}|}\hline
\multicolumn{2}{|c|}{\textbf{List-I}} & \multicolumn{2}{c|}{\textbf{List-II}}\\\hline
(I) & Towards right & (P) & Friction force on P when both blocks are at equilibrium position \& moving right.\\\hline
(II) & Towards left & (Q) & Friction force on Q when both the blocks are at equilibrium position \& moving left.\\\hline
(III) & Zero & (R) & Friction force on P when both the blocks are at right extreme position.\\\hline
(IV) & Maximum & (S) & Friction force on Q when both the blocks are at left extreme position.\\\hline
 & & (T) & Spring force on P when the blocks are moving \& on right side of equilibrium position.\\\hline
\end{tabular}}\par
\begin{qoptions}
\opt{A}{I $\rightarrow$ R,S;II $\rightarrow$ T;III $\rightarrow$ P,Q;IV $\rightarrow$ R,S}
\opt{B}{I $\rightarrow$ R,S;II $\rightarrow$ P,Q;III $\rightarrow$ S,T;IV $\rightarrow$ P,S}
\opt{C}{I $\rightarrow$ R,S;II $\rightarrow$ Q;III $\rightarrow$ R,S,T;IV $\rightarrow$ Q,R}
\opt{D}{I $\rightarrow$ R,S;II $\rightarrow$ Q;III $\rightarrow$ P,R;IV $\rightarrow$ T}
\end{qoptions}
\sect{SECTION-II}
\qn{1} For a partially polarized light, we can define the degree of polarization. The degree of polarization of light is given by $P=\frac{I_{max}-I_{min}}{I_{max}+I_{min}}$. Here $I_{max}$ and $I_{min}$ are maximum and minimum possible intensities in two mutually perpendicular directions. For a given partially polarized light beam, the value of $P$ is 0.25. The approximate ratio of maximum and minimum amplitudes in two mutually perpendicular directions is:
\qn{2} A glass bulb with volumetric expansion coefficient $\alpha=1.50\times10^{-5}\,{}^0\mathrm{C}^{-1}$ is weighed in liquid at temperatures $t=10^\circ$C and $t_1=20^\circ$C. The weights of the displaced liquids are $p=100.01$ g and $p_1=100$ g respectively. Find the volumetric expansion coefficient $\alpha_l$ (in the units of $10^{-5}\,{}^0\mathrm{C}^{-1}$) of the liquid in the temperature interval from $t$ to $t_1$.
\qn{3} A sphere and a solid cube both made of copper are heated up to same temperature and are kept in same surrounding. If radius of sphere and side of cube are same then what will be ratio of initial rate of cooling of sphere to cube ?
\qn{4} Two satellites move in circular orbits around the earth at distances 9000 km and 9010 km from earth's centre. If the satellite which is moving faster has a period of revolution 90 minutes, find the difference in their revolution periods in second ?
\qn{5} An electric field of 300 V/m is confined to a circular area 10 cm in diameter. If the field is increasing at the rate of 20 V/m-s and if the magnitude of magnetic field at a point 15 cm from the centre of the circle is $B\times10^{-18}$ T, find value of $B$.
\qn{6} A uniformly charged sphere is placed inside a charged hollow non-conducting sphere as shown in figure. $O$ is the centre of hollow sphere and $C$ is the centre of solid sphere. The magnitude of charge on both the spheres is $4\left(\sqrt{\frac37}\right)\ \mu$C. The electric field at point `P' which lie just outside the hollow sphere is given by $3\alpha$ kN/C. Find the value of $\alpha$.\par
\qfig{p1_phy_s3_q06_1}
\qn{7} A semicircular disc of radius $R$ and mass $M$ is pulled by a horizontal force $F$ so that is moves with uniform velocity. The coefficient of friction between disc and ground is $\mu=\frac{4}{3\pi}$. If the angle $\theta=\pi/n$ radians, then find the value of $n$.\par
\qfig{p1_phy_s3_q07_1}
\qn{8} A vernier callipers has negative zero error. When both jaws are touching each other, the 7$^{th}$ mark of the vernier coincides with 5$^{th}$ mark on the main scale. In this setup, 10 vernier scale divisions coincides with 9 main scale divisions and length of each main scale division is 1mm. When an object is placed between jaws, main scale reading is at 45$^{th}$ division and vernier coincides is at 5$^{th}$ division mark. What is the true length (in cm) of object ?
\partsec{2}{CHEMISTRY}
\sect{SECTION-I (i)}
\qn{1} For the reaction at 273$^\circ$C\par
N$_2$(g) + 3H$_2$(g) $\longrightarrow$ 2NH$_3$(g)\par
It is observed that the pressure of H$_2$(g) falls from 760 mmHg to 380 mmHg in 100 sec.\par
Select the correct statement(s).
\begin{qoptions}
\opt{A}{Average rate of reaction is $5\times10^{-3}$ atm/sec.}
\opt{B}{Average rate of formation of NH$_3$(g) is $\frac{1}{300}$ atm/sec.}
\opt{C}{Average rate of disappearance of N$_2$(g) is twice of rate of appearance of NH$_3$(g).}
\opt{D}{$\log\left(\frac{-\Delta P_{H_2}}{\Delta t}\right)=\log\left(\frac{-\Delta P_{N_2}}{\Delta t}\right)+0.4771$}
\end{qoptions}
\qn{2} The fixation of N$_2$ by certain bacteria is a biological reduction reaction of great interest. Very few nonbiochemical systems are known which will react with N$_2$ at room temperature and atmospheric pressure. One such compound is the titanium(II) compound titanocene, (C$_5$H$_5$)$_2$Ti sometimes also abbreviated as Cp$_2$Ti, a sandwich compound like ferrocene. Select the correct statement(s) about titanocene (I) \& its dinitrogen complex (II) [Bond length of free N$_2$ 109.8 pm]\par
\qfig{p1_chem_s1_q02_1}
\begin{qoptions}
\opt{A}{The IUPAC name for the complex (II) would be bis(cyclopentadienyl)bis(dinitrogen)titanate(II)}
\opt{B}{NN bond lengths in the complex will be higher than 109.8 pm}
\opt{C}{The NN bond length in the complex is less than 109.8 pm}
\opt{D}{Of the two complexes, only complex (II) follows Sidgwick EAN rule}
\end{qoptions}
\qn{3} Select the correct statement(s) about properties of solid.
\begin{qoptions}
\opt{A}{In antiferromagnetic substance magnetic moments of domains are oppositely oriented and cancel out each other's magnetic moment.}
\opt{B}{H$_2$O, NaCl and C$_6$H$_6$ are weakly repelled by external magnetic field.}
\opt{C}{CrO$_2$ is strongly attracted by external magnetic field.}
\opt{D}{Metals conduct electricity in solid as well as in molten state.}
\end{qoptions}
\qn{4} Which of the following order of stability is correct?
\begin{qoptions}
\opt{A}{\qfigi{p1_chem_s1_q04_A}}
\opt{B}{\qfigi{p1_chem_s1_q04_B}}
\opt{C}{\qfigi{p1_chem_s1_q04_C}}
\opt{D}{\qfigi{p1_chem_s1_q04_D}}
\end{qoptions}
\qn{5} The polymeric structure shown in the diagram contains:\par
\qfig{p1_chem_s1_q05_1}
\begin{qoptions}
\opt{A}{At least one $\alpha$-D-glucopyranose unit}
\opt{B}{Poly-$\beta$-D-fructofuranoside moity}
\opt{C}{(2 $\rightarrow$ 1) $\beta$-glycosidic linkage}
\opt{D}{At least one $\beta$-D-fructofuranose unit}
\end{qoptions}
\qn{6} Which of the following plot(s) is/are correct for a sample of ideal gas in given conditions ?
\begin{qoptions}
\opt{A}{\qfigi{p1_chem_s1_q06_A}}
\opt{B}{\qfigi{p1_chem_s1_q06_B}}
\opt{C}{\qfigi{p1_chem_s1_q06_C}}
\opt{D}{\qfigi{p1_chem_s1_q06_D}}
\end{qoptions}
\sect{SECTION-I (ii)}
\qn{7}\par
{\footnotesize\setlength{\tabcolsep}{2pt}\noindent\begin{tabular}{|c|L{3.3cm}|c|L{3.0cm}|}\hline
\multicolumn{2}{|c|}{\textbf{List-I} (Aqueous solution M = m)} & \multicolumn{2}{c|}{\textbf{List-II}}\\\hline
(P) & Solution-1 : 0.1 M glucose ; Solution-2 : 0.1 M urea & (1) & Solution-1 and Solution-2 are isotonic\\\hline
(Q) & Solution-1 : 0.1 M NaCl ; Solution-2 : 0.1 M Na$_2$SO$_4$ & (2) & Solution-1 has higher freezing point than Solution-2\\\hline
(R) & Solution-1 : 0.1 M BaCl$_2$ ; Solution-2 : 0.1 M urea & (3) & Solution-1 has higher boiling point than Solution-2\\\hline
(S) & Solution-1 : 0.1 M CuSO$_4$ ; Solution-2 : 0.1 M sucrose & (4) & Solution-1 has higher vapour pressure than Solution-2\\\hline
 & & (5) & Solution-1 has greater depression in freezing point than Solution-2\\\hline
\end{tabular}}\par
\begin{qoptions}
\opt{A}{P $\rightarrow$ 3;Q $\rightarrow$ 2,4;R $\rightarrow$ 3,5;S $\rightarrow$ 4,5}
\opt{B}{P $\rightarrow$ 1;Q $\rightarrow$ 3,4;R $\rightarrow$ 2,3;S $\rightarrow$ 3,5}
\opt{C}{P $\rightarrow$ 5;Q $\rightarrow$ 3,5;R $\rightarrow$ 1,5;S $\rightarrow$ 2,5}
\opt{D}{P $\rightarrow$ 1;Q $\rightarrow$ 2,4;R $\rightarrow$ 3,5;S $\rightarrow$ 3}
\end{qoptions}
\qn{8}\par
{\footnotesize\setlength{\tabcolsep}{2pt}\noindent\begin{tabular}{|c|L{2.6cm}|c|L{3.6cm}|}\hline
\multicolumn{2}{|L{3.0cm}|}{\textbf{List-I} (Ionic compound structure)} & \multicolumn{2}{c|}{\textbf{List-II}}\\\hline
(P) & Rock salt structure & (1) & 100\% tetrahedral voids occupied by cation\\\hline
(Q) & Zinc Blende & (2) & 100\% tetrahedral voids occupied by anion\\\hline
(R) & Flourite & (3) & 100\% octahedral voids occupied by cation\\\hline
(S) & Anti flourite (Li$_2$O) & (4) & a = 2r$_+$ + 2r$_-$\\\hline
 & & (5) & $\sqrt3$ a = 4r$_+$ + 4r$_-$\\\hline
\end{tabular}}\par
\begin{qoptions}
\opt{A}{P $\rightarrow$ 3,4;Q $\rightarrow$ 5;R $\rightarrow$ 2,5;S $\rightarrow$ 1,5}
\opt{B}{P $\rightarrow$ 3,4;Q $\rightarrow$ 5;R $\rightarrow$ 1,2;S $\rightarrow$ 1,2}
\opt{C}{P $\rightarrow$ 3,4;Q $\rightarrow$ 2;R $\rightarrow$ 2,5;S $\rightarrow$ 1,5}
\opt{D}{P $\rightarrow$ 1,5;Q $\rightarrow$ 2;R $\rightarrow$ 3,5;S $\rightarrow$ 4}
\end{qoptions}
\qn{9} For the reaction, using suitable alkyl halide and a suitable nucleophile what will be the correct match for the nucleophile used as per List-I and the solvent being used in List-II?\par
{\footnotesize\setlength{\tabcolsep}{2pt}\noindent\begin{tabular}{|c|L{2.3cm}|c|L{3.7cm}|}\hline
\multicolumn{2}{|c|}{\textbf{List-I}} & \multicolumn{2}{c|}{\textbf{List-II}}\\\hline
I & MeS$^{-}$ & P & Reacts with CH$_3$CH$_2$Cl to give SN$^2$ product\\\hline
II & NC$^{-}$ + NH$_4$Cl & Q & Reacts with CH$_3$CHO followed by hydrolysis to give alanine.\\\hline
III & \rule{0pt}{1.9cm}\qfigi[0.75]{p1_chem_s2_q09_1} & R & Reacts with EtX followed by NH$_2$NH$_2$ gives EtNH$_2$\\\hline
IV & Ph$_3$P: + nBuLi & S & Reacts with CH$_3$CH$_2$Cl followed by acetone to give 2-methyl-but-2-ene.\\\hline
 & & T & Reacts with chloro benzene followed by H$_3$O$^+$ give aniline.\\\hline
\end{tabular}}\par
\begin{qoptions}
\opt{A}{I $\rightarrow$ P;II $\rightarrow$ Q;III $\rightarrow$ R;IV $\rightarrow$ S}
\opt{B}{I $\rightarrow$ T;II $\rightarrow$ Q;III $\rightarrow$ P;IV $\rightarrow$ S}
\opt{C}{I $\rightarrow$ T;II $\rightarrow$ S;III $\rightarrow$ R;IV $\rightarrow$ P}
\opt{D}{I $\rightarrow$ P;II $\rightarrow$ S;III $\rightarrow$ T;IV $\rightarrow$ R}
\end{qoptions}
\qn{10}\par
{\footnotesize\setlength{\tabcolsep}{2pt}\noindent\begin{tabular}{|c|L{2.2cm}|c|L{3.9cm}|}\hline
\multicolumn{2}{|c|}{\textbf{List-I} (Species)} & \multicolumn{2}{c|}{\textbf{List-II} (Characteristics of anion)}\\\hline
I & NaN$_3$ & P & Non linear\\\hline
II & KI$_3$ & Q & Linear without using d-orbitals in hybridisation\\\hline
III & KO$_3$ & R & No $\pi$-bond\\\hline
IV & Mg$_2$C$_3$ & S & Paramagnetic\\\hline
 & & T & Hydrolysed product have 6$\sigma$ and 2$\pi$ bonds.\\\hline
\end{tabular}}\par
\begin{qoptions}
\opt{A}{I $\rightarrow$ Q;II $\rightarrow$ R;III $\rightarrow$ R;IV $\rightarrow$ P}
\opt{B}{I $\rightarrow$ Q;II $\rightarrow$ R;III $\rightarrow$ S;IV $\rightarrow$ T}
\opt{C}{I $\rightarrow$ P;II $\rightarrow$ Q;III $\rightarrow$ S;IV $\rightarrow$ T}
\opt{D}{I $\rightarrow$ S;II $\rightarrow$ P;III $\rightarrow$ Q;IV $\rightarrow$ R}
\end{qoptions}
\sect{SECTION-II}
\qn{1} Number of halide ions in formula of cryolite $=x$\par
No. of metal atoms in the formula of chromite $=y$\par
Oxidation number of metal in cassiterite $=z$\par
Find $(100x+50y+z)$
\qn{2} C$_4$H$_8$O will have $x$ open chain carbonyl compounds, $y$ unsaturated alcohols and $z$ ethers ($z_1$ - aliphatic, $z_2$ - alicyclic). What is $[100x+10y+z-0.1\,(z_1-z_2)]$ ?
\qn{3} A second period element X forms of 5 different oxidation state binary oxides some of these are neutral, mostly are acidic. Element Y is the congener (the very next member of the same group) of X.\par
X forms only one chloride (I), Y forms two (II \& III). If one mole each of I, II \& III is hydrolysed in enough water. Find the total number of HCl molecules formed by hydrolysis:
\qn{4} H$_4$EDTA combines with calcium in basic medium to form a very stable complex. This complex has $x$ membered rings ; $y$ number of such rings are present in the complex. The number of stereoisomers of the complex is $z$. Find $[100(x)+10(y)+z]$
\qn{5} In the structure of tetranuclear anion of borax, total number of B$-$O$-$B bonds are $x$, maximum number of boron atoms in same plane are $y$ and total number of lone pairs are $z$.\par
Find $(100x+10y+z)$
\qn{6} Find sum of unpaired electron(s) in :\par
(i) Atomic nitrogen\par
(ii) Unstable S$_2$ molecule\par
(iii) Potasium superoxide\par
(iv) [Ni(en)$_3$]SO$_4$
\qn{7} In an ore the only oxidisable material is Sn$^{2+}$. This ore is treated with a dichromate solution containing 4.9 g K$_2$Cr$_2$O$_7$ in 1 litre. A 0.50 g sample of the ore required 20 ml of the titrant to reach equivalent point. Find mass percent of Sn in ore.\par
[Atomic mass of Sn = 119, K = 39, Cr = 52]
\qn{8} The conductivity of a saturated solution of Ag$_3$PO$_4$ is $9\times10^{-6}$ S m$^{-1}$ and its equivalent conductivity is $1.50\times10^{-4}$ Sm$^2$ equivalent$^{-1}$. The K$_{sp}$ of Ag$_3$PO$_4$ is $x\times10^{-18}$, then value of $x$ is :
\partsec{3}{MATHEMATICS}
\sect{SECTION-I (i)}
\qn{1} Which of the following statement is/are correct ?
\begin{qoptions}
\opt{A}{Tangents at $P$, $Q$ on the parabola $y^2=4ax$ intersect at $T$. If $S$ be the focus of parabola then $ST^2=SP\times SQ$}
\opt{B}{circle passing through co-normal points also passes through vertex of parabola}
\opt{C}{The ortho-centre of triangle formed by point of intersection of Tangents of parabola lies on directrix}
\opt{D}{The Quadrilateral $SPMT$ where $S$ is focus, $P$ is any point on the parabola, $M$ is foot of perpendicular from $P$ to directrix and $T$ is point of intersection of tangent at $P$ with axis of parabola is rhombus.}
\end{qoptions}
\qn{2} Identify the correct statements
\begin{qoptions}
\opt{A}{If $f(x)>0$, $f''(x)>0\ \forall\ x\in[0,1]$ and $f(0)=0$, $f(1)=1$, then $f(x)\cdot f^{-1}(x)<x^2\ \forall\ x\in(0,1)$}
\opt{B}{If $f(x)$ is a twice derivable function and $f(\alpha)=f(\beta)=0$, $(\alpha<\beta)$ then $(f'(x))^2+f(x)f''(x)=0$ has at least two real roots in $x\in(\alpha,\beta)$}
\opt{C}{Let $g:[a,b]\rightarrow R$ be twice differentiable such that $g''(x)+g'(x)\cdot g(x)=g(x)\ \forall\ x\in[a,b]$ and $g(a)=g(b)=0$, then $g(x)=0\ \forall\ x\in[a,b]$}
\opt{D}{Let $f:[-1,1]\rightarrow R$ be a differentiable function such that $f(0)=0$ and $f(x)>0\ \forall\ x\in(-1,1)-\{0\}$, then $\frac{f'(c)}{f(c)}=\frac32\frac{f'(1-3c)}{f(1-3c)}+\frac{1}{1-2c}$ for atleast one $c\in\left(\frac13,\frac12\right)$}
\end{qoptions}
\qn{3} Let $F$ be the foot of perpendicular from $P(1,2,-3)$ on the line $L:\frac{x+1}{2}=\frac{y-3}{-2}=\frac{z}{-1}$. Given $Q(x_1,y_1,z_1)$ and $R(x_2,y_2,z_2)$ are points at a distance of 3 units from $F$ on line $L$. $\pi_1$ and $\pi_2$ be two planes passing through $Q$ and $R$ respectively and perpendicular to the line of intersection to the planes $9x-7y+6z+48=0$, $x+y-z=7$. Which of the following statement is/are correct?
\begin{qoptions}
\opt{A}{The coordinates of $F$ is $(1,1,-1)$}
\opt{B}{$\sum_{i=1}^{2}(x_i^2+y_i^2+z_i^2)=24$}
\opt{C}{The distance between $\pi_1$ and $\pi_2=\frac{81}{\sqrt{182}}$}
\opt{D}{The normal vector of plane $\pi_1$ and $\pi_2$ is parallel to the vector $\hat i-15\hat j+16\hat k$}
\end{qoptions}
\qn{4} Let $S$ be the circle passes through foci of an ellipse $(E)$ and the ellipse $(E)$ intersect the hyperbola $(H)$ $2x^2-2y^2=1$ orthogonally. The eccentricity of the ellipse is half of that of hyperbola. If the axes of the ellipse are along the coordinate axes, $S$ is also passes through focus of the parabola $(y-3)^2=12(x+1)$ then
\begin{qoptions}
\opt{A}{Radius of director circle of $S_1$ is $\sqrt{10}$}
\opt{B}{Equation of ellipse is $x^2+2y^2=2$}
\opt{C}{Length of tangent from $(-5,1)$ to circle is $\sqrt{21}$}
\opt{D}{Equation of ellipse is $2x^2+y^2=4$}
\end{qoptions}
\qn{5} $I_n=\displaystyle\int_{n/2}^{\frac{n+1}{2}}\frac{\sin(\pi\sin^2\pi x)}{(2)^{\frac x2}}\,dx,\ n\in I$
\begin{qoptions}
\opt{A}{$\frac{I_{2022}}{I_{2026}}=2$}
\opt{B}{$\frac{I_{2022}}{I_{2026}}=\frac{1}{\sqrt2}$}
\opt{C}{$\frac{\sum_{n=0}^{\infty}I_{8n}}{I_0}=\frac43$}
\opt{D}{$\frac{\sum_{n=0}^{\infty}I_n}{I_0}=2$}
\end{qoptions}
\qn{6} Which of the following statement is correct?
\begin{qoptions}
\opt{A}{$30!>3!5!10!15!$}
\opt{B}{$4\sin^{-1}\left(\frac14\right)+3\sin^{-1}\left(\frac13\right)>7\sin^{-1}\left(\frac27\right)$}
\opt{C}{Consider a real system of equations $x^2+y^2=4$, $|x|+|y|=k$. Then the system has always even number of solutions.}
\opt{D}{$f(x)=\begin{cases}1-|x|, & |x|\le1\\ 0, & |x|>1\end{cases}$ and $g(x)=f(x-1)+f(x+1)$\par Then the number of points where $g(x)$ is not differentiable is 4.}
\end{qoptions}
\sect{SECTION-I (ii)}
\qn{7} Match the following\par
{\footnotesize\setlength{\tabcolsep}{2pt}\noindent\begin{tabular}{|c|L{4.6cm}|c|L{1.6cm}|}\hline
\multicolumn{2}{|c|}{\textbf{List-I}} & \multicolumn{2}{c|}{\textbf{List-II}}\\\hline
I. & Domain of the function\par\resizebox{\linewidth}{!}{$f(x)=\log_{|x-1|-|x+1|-\frac12}\left[\left(x^2-x+1\right)^{x^2-3x+2}-1\right]$}\par & P. & $[10,\infty)$\\\hline
II. & Consider the curve $(2x+y-1)^2=5(x-2y-3)$ then the possible coordinates of foci & Q. & $\left(\frac34,-\frac12\right)$\\\hline
III. & Let $f(x)=\begin{cases}2-x+a^2-9a-9, & x<2\\ 2x-3, & x\ge2\end{cases}$. Where $a>0$. If $f(x)$ has local minima at $x=2$, then $a\in$ & R. & $\left(-\infty,-\frac14\right)-\left\{-\frac34\right\}$\\\hline
IV. & If $\cos^2x-(a-1)\cos x+2a\ge6\ \forall\ x\in R$, then the true set of values of $a$ is & S. & $[4,\infty)$\\\hline
 & & T. & $\left[-\frac13,5\right)$\\\hline
\end{tabular}}\par
\begin{qoptions}
\opt{A}{I $\rightarrow$ R;II $\rightarrow$ Q;III $\rightarrow$ P;IV $\rightarrow$ S}
\opt{B}{I $\rightarrow$ Q;II $\rightarrow$ R;III $\rightarrow$ P;IV $\rightarrow$ S}
\opt{C}{I $\rightarrow$ S;II $\rightarrow$ R;III $\rightarrow$ P;IV $\rightarrow$ Q}
\opt{D}{I $\rightarrow$ S;II $\rightarrow$ Q;III $\rightarrow$ P;IV $\rightarrow$ R}
\end{qoptions}
\qn{8} Match the following\par
{\footnotesize\setlength{\tabcolsep}{2pt}\noindent\begin{tabular}{|c|L{4.8cm}|c|L{1.2cm}|}\hline
\multicolumn{2}{|c|}{\textbf{List-I}} & \multicolumn{2}{c|}{\textbf{List-II}}\\\hline
I. & $S=\{8,9,10,\ldots\ldots\}$ $f:S\rightarrow S$. For $x,y\in S$ and $f(x+y)=f(xy)\ \forall x,y\ge4$; $f(8)=9$ then $2f(9)=$ & P. & 1\\\hline
II. & $x,y\in R^+$ such that $x+y=2$ then the greatest value of $x^3y^3(x^3+y^3)$ is & Q. & 98\\\hline
III. & If\par\resizebox{\linewidth}{!}{$f(x)=\left(\left[\{x\}\right]\tan^{-1}\left(\frac{x^2-3x-1}{x^2-3x+5}\right)+3-x^7\right)^{1/7}$}\par. $[x]$ is G.I.F, $\{x\}$ fractional part then $f^{-1}(50)-f(50)+f(f(100))$ & R. & 18\\\hline
IV. & \resizebox{\linewidth}{!}{$\left|\frac{\left(1+\sqrt3\tan1^\circ\right)\left(1-\sqrt3\tan1^\circ\right)}{\left(\sqrt3+\tan1^\circ\right)\left(\sqrt3-\tan1^\circ\right)}\left(\frac{\tan89^\circ}{\cot3^\circ}\right)\right|$} & S. & 2\\\hline
 & & T. & 100\\\hline
\end{tabular}}\par
\begin{qoptions}
\opt{A}{I $\rightarrow$ R;II $\rightarrow$ S;III $\rightarrow$ P;IV $\rightarrow$ Q}
\opt{B}{I $\rightarrow$ P;II $\rightarrow$ S;III $\rightarrow$ R;IV $\rightarrow$ T}
\opt{C}{I $\rightarrow$ T;II $\rightarrow$ S;III $\rightarrow$ R;IV $\rightarrow$ P}
\opt{D}{I $\rightarrow$ R;II $\rightarrow$ S;III $\rightarrow$ T;IV $\rightarrow$ P}
\end{qoptions}
\qn{9} Match the following\par
{\footnotesize\setlength{\tabcolsep}{2pt}\noindent\begin{tabular}{|c|L{4.8cm}|c|L{1.2cm}|}\hline
\multicolumn{2}{|c|}{\textbf{List-I}} & \multicolumn{2}{c|}{\textbf{List-II}}\\\hline
I. & The number of integral solutions of the equation $4\displaystyle\int_0^\infty\frac{\ell nt}{x^2+t^2}dt=\pi\ell n2,\ t>0$ & P. & 12\\\hline
II. & $\displaystyle\int_1^e\frac{x^4\ell nx+2}{x^3\ell nx+x}dx=\frac{e^2+a}{b}-\ell n\left(e^2+1\right)$, $a,b\in N$. Then $a+b=$ & Q. & 0\\\hline
III. & The number of solutions of the equation $2\tan^{-1}\left(\sec^2\pi x\right)=\sin^{-1}\left(x^3-x^2+x+2\right)$ is & R. & 5\\\hline
IV. & $f:R\rightarrow R$, $f(x)=x^4-x^2+1$. The number of values of $x$ for which $f(f(f(x)))\le x^8$ is $K$ then $K+10=$ & S. & 2\\\hline
 & & T. & 15\\\hline
\end{tabular}}\par
\begin{qoptions}
\opt{A}{I $\rightarrow$ T;II $\rightarrow$ S;III $\rightarrow$ Q;IV $\rightarrow$ P}
\opt{B}{I $\rightarrow$ S;II $\rightarrow$ R;III $\rightarrow$ S;IV $\rightarrow$ Q}
\opt{C}{I $\rightarrow$ S;II $\rightarrow$ R;III $\rightarrow$ Q;IV $\rightarrow$ P}
\opt{D}{I $\rightarrow$ P;II $\rightarrow$ Q;III $\rightarrow$ T;IV $\rightarrow$ P}
\end{qoptions}
\qn{10} Match the following\par
{\footnotesize\setlength{\tabcolsep}{2pt}\noindent\begin{tabular}{|c|L{4.8cm}|c|L{1.2cm}|}\hline
\multicolumn{2}{|c|}{\textbf{List-I}} & \multicolumn{2}{c|}{\textbf{List-II}}\\\hline
I & If the foot of perpendicular drawn from $(1,9,7)$ to the line passing through the point $(3,2,1)$ and parallel to the planes $x+2y+z=0$ and $3y-z=3$ is $(p,q,r)$ then $2p+q+r$ is equal to & P & 1\\\hline
II & Consider $f(x)=\sin2\pi x+\{x\}$, $x\in[0,10]$, number of points where $f$ attains local maximum is (Where $\{.\}$ represents fractional part) & Q & 10\\\hline
III & If $\displaystyle\int_0^{\pi/4}x\cdot\left(\prod_{n=1}^{\infty}\cos\left(\frac{x}{2^n}\right)\right)dx=\frac{p-\sqrt q}{2}$ $(p,q\in N)$ then $p+q=$ & R & 3\\\hline
IV & $\displaystyle\int_0^e x\sqrt{x\cdot\sqrt[3]{x\cdot\sqrt[4]{x\cdot\sqrt[5]{x\ldots\infty}}}}\,dx=e^{A(e)+B}$, $A,B\in$ Integer, then $A+B=$ & S & 0\\\hline
 & & T & 4\\\hline
\end{tabular}}\par
\begin{qoptions}
\opt{A}{I $\rightarrow$ R;II $\rightarrow$ Q;III $\rightarrow$ T;IV $\rightarrow$ S}
\opt{B}{I $\rightarrow$ R;II $\rightarrow$ P;III $\rightarrow$ T;IV $\rightarrow$ Q}
\opt{C}{I $\rightarrow$ P;II $\rightarrow$ R;III $\rightarrow$ T;IV $\rightarrow$ Q}
\opt{D}{I $\rightarrow$ S;II $\rightarrow$ P;III $\rightarrow$ Q;IV $\rightarrow$ T}
\end{qoptions}
\sect{SECTION-II}
\qn{1} $\frac{((3!)!)!}{6!}=n!$, then number of zero's at the end in $n!$
\qn{2} Given that $\log_2\sin x+\log_2\cos x=-1$ and $\log_2(\sin x+\cos x)=\frac12(\log_{10}n-1)$, find $n$ ?
\qn{3} In the expansion of $(ax+b)^{2000}$ where $a$ and $b$ are relatively prime, the coefficients of $x^2$ and $x^3$ are equal, find $a+b=$ \underline{\hspace{1cm}}
\qn{4} The equation $2000x^6+100x^5+10x^3+x-2=0$ has exactly two real roots, one of which is $\frac{m+\sqrt n}{r}$, where $m$ is nonzero integer and $n$ and $r$ are relatively prime natural numbers, find $(m+n+r)$.
\qn{5} Consider the polynomials $P(x)=x^6-x^5-x^3-x^2-x$ and $Q(x)=x^4-x^3-x^2-1$. Given that $z_1,z_2,z_3,z_4$ are roots of $Q(x)=0$, find $P(z_1)+P(z_2)+P(z_3)+P(z_4)$
\qn{6} Let $a=\frac{\pi}{2008}$. Find the smallest positive integer $n$ such that\par
$2\Big[\cos a\sin a+\cos(4a)\cdot\sin(2a)+\cos(9a)\sin(3a)$\par$+\ldots.+\cos\left(n^2a\right)\cdot\sin(na)\Big]$\par
is an integer.
\qn{7} Consider three matrices $A=[a_{ij}]_{3\times5}$, $B=[b_{ij}]_{5\times10}$ and $C=AB=[c_{ij}]_{3\times10}$. It is given that $a_{ij}=\frac{i}{3j-2}$, $b_{ij}=\frac{j}{3i+1}$, then the value of the element $C_{28}$ is \underline{\hspace{1.5cm}}.
\qn{8} The differential equation $\frac{dy}{dx}=\frac{x\sqrt{x^2-1}+y}{\sqrt{x^2-1}}$ which gives $y(1)=1$, then $|y(-1)|=$ \underline{\hspace{1cm}} ($|.|$ represent modulus)
\raggedcolumns
\papersec{2}
\partsec{1}{PHYSICS}
\sect{SECTION-I (i)}
\qn{1} A liquid is flowing in a horizontal pipe of uniform cross section at a speed $v$. Two tubes A and B are inserted into the pipe as shown. Assume the flow to remain streamline inside the pipe. The difference in height of the liquid in two tubes is :-\par
\qfig{p2_phy_s1_q01_1}
\begin{qoptions}
\opt{A}{$\frac{v^2}{g}$}
\opt{B}{$\frac{v^2}{2g}$}
\opt{C}{$\frac{2v^2}{g}$}
\opt{D}{Information is not enough to calculate.}
\end{qoptions}
\qn{2} As shown in the figure, a particle A of mass $2m$ and carrying charge $q$ is connected by a light rigid rod of length $L$ to another particle B of mass $m$ and carrying charge $-q$. The system is placed in an electric field $\vec E$. The rod can oscillate in horizontal plane. The electric force on a charge $q$ in an electric field $\vec E$ is $\vec F=q\vec E$. After the system settles into equilibrium, one particle is given a small push in the transverse direction so that the rod makes a small angle $\theta_0$ with the electric field. Find maximum tension in the rod.\par
\qfig{p2_phy_s1_q02_1}
\begin{qoptions}
\opt{A}{$qE+qE\theta_0^{\,2}$}
\opt{B}{$qE+\frac{qE\theta_0^2}{4}$}
\opt{C}{$qE+\frac{qE\theta_0^2}{3}$}
\opt{D}{$qE+\frac{qE\theta_0^2}{6}$}
\end{qoptions}
\qn{3} In the given circuit switch S is closed at $t=0$. The current $I$ in the figure at time $t$ is :-\par
\qfig{p2_phy_s1_q03_1}
\begin{qoptions}
\opt{A}{$i=\frac{\varepsilon e^{-t/RC}}{2R}$}
\opt{B}{$i=\frac{2\varepsilon e^{-t/RC}}{3R}$}
\opt{C}{$i=\frac{\varepsilon e^{-\frac{2t}{3RC}}}{2R}$}
\opt{D}{$i=\frac{\varepsilon e^{-t/RC}}{R}$}
\end{qoptions}
\qn{4} A cylindrical resistance is connected across battery $\varepsilon$. Cylinder has uniform free electron density, mid part of cylinder has larger radius as shown in figure. Then graph for $V_d$ (drift velocity) V/S (distance across the length of the resistance) $L$ will be :-\par
\qfig{p2_phy_s1_q04_1}
\begin{qoptions}
\opt{A}{\qfigi{p2_phy_s1_q04_A}}
\opt{B}{\qfigi{p2_phy_s1_q04_B}}
\opt{C}{\qfigi{p2_phy_s1_q04_C}}
\opt{D}{\qfigi{p2_phy_s1_q04_D}}
\end{qoptions}
\sect{SECTION-I (ii)}
\qn{5} In the given figure if A is connected with C and B is connected with D, then\par
\qfig{p2_phy_s2_q05_1}
\begin{qoptions}
\opt{A}{Charge flows in the circuit is 12 $\mu$C.}
\opt{B}{Heat produced in the circuit is 60 $\mu$J.}
\opt{C}{Charge flows in the circuit is 14 $\mu$C.}
\opt{D}{Heat produced in the circuit is 490 $\mu$J.}
\end{qoptions}
\qn{6} When an AC source having frequency $\frac{50}{\pi}$ Hz applied in series with a circuit having resistance $R=\frac{1}{\sqrt3}$ k$\Omega$ and inductor or capacitor then $i$ and $E$ as a function of time are observed to vary with time as shown in figure and phase difference between them is $60^\circ$.\par
\qfig{p2_phy_s2_q06_1}
\begin{qoptions}
\opt{A}{Circuit has inductor in series with resistor}
\opt{B}{Circuit has capacitor in series with resistor}
\opt{C}{Value of $L$ or $C$ is $(10\times10^{-6})$ SI unit}
\opt{D}{Value of $L$ or $C$ is $(3\times10^{-6})$ SI unit}
\end{qoptions}
\qn{7} A time varying magnetic field exists in the region, $\vec B=B_0t\left(\hat i+\hat k\right)$. $ABC$ is a conducting triangular loop (see figure). Then:\par
\qfig{p2_phy_s2_q07_1}
\begin{qoptions}
\opt{A}{Direction of induced current in the loop is from $A$ to $B$.}
\opt{B}{Direction of induced current in the loop is from $B$ to $A$.}
\opt{C}{Induced EMF in loop is $B_0a^2$.}
\opt{D}{Induced EMF in loop is $\sqrt2B_0a^2$.}
\end{qoptions}
\qn{8} A screw gauge (as shown in figure) having 100 equal division on circular scale and a pitch of 1 mm is used to measure the diameter of a wire of length 4.2 cm. The main scale reading is 2 mm and 45$^{th}$ circular division coincides with the main scale. $\left(\pi=\frac{22}{7}\right)$. Choose the correct option/options :\par
\qfig{p2_phy_s2_q08_1}
[Before the measurement when screw is taut]
\begin{qoptions}
\opt{A}{Diameter of wire is 2.50 mm}
\opt{B}{Diameter of wire is 2.40 mm}
\opt{C}{Surface area (in cm$^2$) of wire in appropriate significant figure is 3.2.}
\opt{D}{Surface area (in cm$^2$) of wire in appropriate significant figure is 3.168}
\end{qoptions}
\qn{9} The x-ray spectrum of a metallic target has been shown in figure. Given : $R=1.1\times10^7\mathrm{m}^{-1}$ and $hc=1240$eV$-$nm [Both characteristic x-rays belong to k series ($k_\alpha$ \& $k_\beta$)]. Choose the correct option/options :\par
\qfig{p2_phy_s2_q09_1}
\begin{qoptions}
\opt{A}{wave length of k$_\alpha$ x-ray is 0.286 nm}
\opt{B}{wavelength of k$_\alpha$ x-ray is 0.303 nm}
\opt{C}{atomic number of the target atom is 21.}
\opt{D}{accelerating voltage for x-rays tube is 12400 V.}
\end{qoptions}
\qn{10} A narrow capillary tube open at both ends of radius $r=0.25$ mm contains some water in it. When the tube is held vertical, the water remains inside the tube as shown. Choose the correct option/options. [Given : surface tension $=0.075$ N/m and contact angle $=0^\circ$, $g=10\frac{\mathrm m}{\mathrm s^2}$]\par
\qfig{p2_phy_s2_q10_1}
\begin{qoptions}
\opt{A}{For $h=4$cm, lower meniscus is convex upwards.}
\opt{B}{For $h=6$ cm, lower meniscus is flat.}
\opt{C}{For $h=8$ cm, lower menicus is concave up.}
\opt{D}{For all values of $h$ lower meniscus is convex upward always.}
\end{qoptions}
\sect{SECTION-III}
\qn{1} A uniform disc of radius $R=0.2$ m is projected along a rough horizontal surface with initial linear velocity $v_0=2$m/s and initial angular velocty $\omega_0=5$ rad/sec as shown in the figure. Coefficient of friction is $\mu=0.1$. The time (in second) at which it starts pure rolling is.\par
\qfig{p2_phy_s3_q01_1}
\qn{2} A hydrogen atom in excited state $n$ reaches the ground state by emitting three photons of frequency in ratio 5 : 72 : 243. Wavelength corresponding to one of the transitions belongs to visible region. Find the value of `n'.
\qn{3} In the fusion reaction ${}^2_1$H $+$ ${}^2_1$H $\rightarrow$ ${}^3_2$He $+$ ${}^1_0$n, the masses of deuteron, helium and neutron expressed in a.m.u. are 2.015, 3.017 and 1.009 respectively. If 1 kg of deuterium undergoes complete fusion, the amount of total energy released is $a\times10^{13}$ (in Joules), find nearest integer of $a$. (Given 1 a.m.u.$\equiv931.5$ MeV/c$^2$)
\qn{4} Two moles of a diatomic gas undergoes two process in succession. Starting from temperature $T_1$, it is first heated at constant volume until its pressure is tripled. It is further heated at constant pressure until its volume is doubled. What is the average molar heat capacity for the whole process? If your answer is $nR$, where $R$ is universal gas constant. Round off $n$ to nearest integer.
\qn{5} A stone of mass 0.2 kg is tied to one end of a thread of length 0.1 m whirled in a vertical circle. When the stone is at the lowest point of circle, tension in thread is 52N, then velocity (in m/s) of the stone will be :$-$
\qn{6} In the circuit shown in figure, calculate the current (in A)through S in the circuit when S is closed.\par
\qfig{p2_phy_s3_q06_1}
\qn{7} One end of a string of length $L$ is tied to the ceiling of a lift accelerating upwards with an acceleration $2g$. The other end of the string is free. The linear mass density of the string varies linearly from 0 to $\lambda$ from bottom to top. The acceleration of the wave w.r.t. string through out string is $\frac{pg}{4}$. Find $p$.
\qn{8} Suppose the potential energy between an electron and a proton at a distance $r$ is $-\frac{Ke}{3r^3}$. The velocity of orbital motion of the electron is given as, $\frac{\alpha}{\beta}\left(\frac{n^3h^3}{\pi^3m^2ke^2}\right)$. Then find the value of $(\alpha+\beta)$.
\raggedcolumns
\partsec{2}{CHEMISTRY}
\sect{SECTION-I (i)}
\qn{1} Which of the following does not liberate NO$_2$ gas
\begin{qoptions}
\opt{A}{AgNO$_3$ $\xrightarrow{\Delta}$}
\opt{B}{NaNO$_3$ $\xrightarrow[500^0\mathrm C]{\Delta\ \mathrm{below}}$}
\opt{C}{Hg(NO$_3$)$_2$ $\xrightarrow{\Delta}$}
\opt{D}{Cu(NO$_3$)$_2$ $\xrightarrow{\Delta}$}
\end{qoptions}
\qn{2} During the Hinsberg Test, which of the following primary amines is most likely to be detected as a secondary amine?
\begin{qoptions}
\opt{A}{MeNH$_2$}
\opt{B}{\qfigi{p2_chem_s1_q02_B}}
\opt{C}{\qfigi{p2_chem_s1_q02_C}}
\opt{D}{\qfigi{p2_chem_s1_q02_D}}
\end{qoptions}
\qn{3} \qfigi{p2_chem_s1_q03_1}\par
Find the product `X'
\begin{qoptions}
\opt{A}{\qfigi{p2_chem_s1_q03_A}}
\opt{B}{\qfigi{p2_chem_s1_q03_B}}
\opt{C}{\qfigi{p2_chem_s1_q03_C}}
\opt{D}{\qfigi{p2_chem_s1_q03_D}}
\end{qoptions}
\qn{4} Calculate change in Gibbs energy of 10 moles of ice when pressure on it is increased from 1 bar to 101 bar isothermally at $-73^\circ$C.\par
[Given : Density of ice at $-73^\circ$C $=0.9$ g mL$^{-1}$]
\begin{qoptions}
\opt{A}{1000 J}
\opt{B}{2000 J}
\opt{C}{76360 J}
\opt{D}{4000 J}
\end{qoptions}
\sect{SECTION-I (ii)}
\qn{5} Which of the following molecular orbitals has two nodal planes?
\begin{qoptions}
\opt{A}{$\pi_{2P_y}$}
\opt{B}{$\pi^*_{2P_x}$}
\opt{C}{$\sigma^*_{2P_z}$}
\opt{D}{$\pi^*_{2P_y}$}
\end{qoptions}
\qn{6} The compound Na$_2$IrCl$_6$ reacts with triphenylphosphine in diethyleneglycol in an atmosphere of CO to give $\left[\mathrm{IrCl(CO)(PPh_3)_2}\right]$, known as Vaskas compound. (Atomic number of Ir $=77$)\par
Select correct for Vaskas compound :
\begin{qoptions}
\opt{A}{The IUPAC name of the complex is carbonylchloridobis (triphenylphosphine)iridium(I)}
\opt{B}{The hybridization of the metal ion is sp$^3$}
\opt{C}{The magnetic moment (spin only) of the complex is zero}
\opt{D}{The complex shows geometrical as well as optical isomerism}
\end{qoptions}
\qn{7} Which of the following is incorrect when alkali metal is allowed to react with the increased concentrated in liq.NH$_3$
\begin{qoptions}
\opt{A}{Paramagnetic character of solvated electrons gets changed}
\opt{B}{Solvated electrons associate to form electron pairs}
\opt{C}{Oxidising nature of the solution is increased}
\opt{D}{Color of solution remain same as in dilute}
\end{qoptions}
\qn{8} The major product formed in the following reaction is\par
\qfig{p2_chem_s2_q08_1}
\begin{qoptions}
\opt{A}{\qfigi{p2_chem_s2_q08_A}}
\opt{B}{\qfigi{p2_chem_s2_q08_B}}
\opt{C}{\qfigi{p2_chem_s2_q08_C}}
\opt{D}{\qfigi{p2_chem_s2_q08_D}}
\end{qoptions}
\qfullbegin
\qn{9}\par
\qfig{p2_chem_s2_q09_1}
\qfullsplit
\begin{qoptions}
\opt{A}{\qfigi{p2_chem_s2_q09_A}}
\opt{B}{\qfigi{p2_chem_s2_q09_B}}
\opt{C}{\qfigi{p2_chem_s2_q09_C}}
\opt{D}{\qfigi{p2_chem_s2_q09_D}}
\end{qoptions}
\qfullend
\qn{10} At 25$^0$C, for half-cell reactions : M + I$^{-}$ $\rightarrow$ MI + e$^{-}$, E$^0=0.152$V\par
M $\rightarrow$ M$^+$ + e$^{-}$, E$^0=-0.80$V , M is a solid metal.\par
Then the value of log K$_{sp}$ for the sparingly soluble salt MI is ?
\begin{qoptions}
\opt{A}{$-16.13$}
\opt{B}{$-20.22$}
\opt{C}{$-8.05$}
\opt{D}{$-32.26$}
\end{qoptions}
\sect{SECTION-III}
\qn{1} A reaction follows temperature - rate constant relation as\par
$\ell nk=30-\frac{500}{T}$\par
If $R=2$ cal k$^{-1}$ mol$^{-1}$ then the value of activation energy of reaction (in kcal/mole) is :
\qn{2} A first order chemical reaction.A $\rightarrow$ P is occurring at 500 K. If the same reaction is to be carried out in presence of a catalyst, temperature of 400 K is required to be maintained in order to maintain the same rate of reaction. If the catalyst lowers the activation energy by 30 kJ, then the value of activation energy of the original pathway (uncatalysed reaction) is $x$ kJ/mol, then $\frac{x}{25}$ is
\qn{3} A closed container of volume 0.02 m$^3$ contains a mixture of Ne and Ar at a temperature of 27$^0$C and a pressure of $10^5$N/m$^2$. If the total mass of mixture is 28 gm. Find the mass of Neon in grams?\par
[Given : Atomic mass : Ne = 20, Ar = 40 and R = 25/3 J mol$^{-1}$ K$^{-1}$]
\qn{4} The enthalpy of fusion of ice is 6.02kJ mol$^{-1}$. The heat capacity of water is 4.214 J g$^{-1}$ K$^{-1}$. What is the smallest number of ice cubes at 0$^\circ$ C, each containing one mole of H$_2$O , that are needed to cool 500g of liquid water from 20$^\mathrm{o}$C to 0$^\mathrm{o}$C ?
\qn{5} How many of of the following species are aromatic\par
\qfig{p2_chem_s3_q05_1}
\qfig{p2_chem_s3_q05_2}
\qn{6} When mixture of\par\qfig{p2_chem_s3_q06_1}\par is reacted with Br$_2$ / OH$^{-}$, how many possible amines are obtained?
\qn{7} The number of only acidic oxides among the following oxides is\par
CO, CO$_2$, SiO$_2$, PbO, PbO$_2$, SnO, NO, N$_2$O, N$_2$O$_3$, NO$_2$, P$_4$O$_6$, P$_4$O$_8$, P$_4$O$_9$, P$_4$O$_{10}$
\qn{8} The number of compounds among the following compounds which give characteristic colour or precipitate with neutral ferric chloride solution is\par
\qfig{p2_chem_s3_q08_1}
\raggedcolumns
\partsec{3}{MATHEMATICS}
\sect{SECTION-I (i)}
\qn{1} Let $C_1:x^2+y^2=1$, $C_2:(x-10)^2+y^2=1$ and $C_3:x^2+y^2-10x-42y+457=0$ are three circles. A circle $C$ has been drawn to touch circle $C_1$ and $C_2$ externally \& $C_3$ internally, Now circle $C_1$, $C_2$ \& $C_3$ start rolling on circumference of $C$ in anticlockwise direction with constant speed. The locus of the centroid of triangle formed by the centres of $C_1$, $C_2$ \& $C_3$ as its vertices -
\begin{qoptions}
\opt{A}{$x^2+y^2-12x-22y+144=0$}
\opt{B}{$x^2+y^2-10x-24y+144=0$}
\opt{C}{$x^2+y^2-8x-20y+64=0$}
\opt{D}{$x^2+y^2-4x-2y-4=0$}
\end{qoptions}
\qn{2} Define the function $A,L,E,N$ as follows\par
$A(x)=24x$\par
$L(x)=12x^2+12x$\par
$E(x)=4x^3+12x^2+8x$\par
$N(x)=x^4+6x^3+11x^2+6x$\par
The value of $\displaystyle\sum_{K=1}^{10}\left(\frac{A(K)+L(K)+E(K)+N(K)}{{}^{11}C_2+{}^{12}C_3+{}^{13}C_4+{}^{14}C_5}\right)$ is
\begin{qoptions}
\opt{A}{24}
\opt{B}{25}
\opt{C}{28}
\opt{D}{34}
\end{qoptions}
\qn{3} If $f(x)$ is continuous, increasing and an odd function such that $\displaystyle\int_{-1}^{4}f(x)\,dx=10$ and $\displaystyle\int_0^1f(x)\,dx=\frac32$ then the area bounded by $y=f(x)$, $x$-axis in between the ordinates $x=-4$ and $x=4$ is
\begin{qoptions}
\opt{A}{23}
\opt{B}{19}
\opt{C}{20}
\opt{D}{$\frac{23}{2}$}
\end{qoptions}
\qn{4} The number of shortest paths from point $A$ to $D$\par
(as shown in figure)\par
\qfig{p2_math_s1_q04_1}
\begin{qoptions}
\opt{A}{276}
\opt{B}{186}
\opt{C}{150}
\opt{D}{126}
\end{qoptions}
\sect{SECTION-I (ii)}
\qn{5} Let $z$ be complex number such that $|z|=1$ then possible values of $|1-z|+|1+z^2|$ among the following is/are
\begin{qoptions}
\opt{A}{$\frac{1}{\sqrt2}$}
\opt{B}{$\sqrt2$}
\opt{C}{$\pi$}
\opt{D}{4}
\end{qoptions}
\qn{6} Let $\alpha,\beta$ $(\alpha>\beta)$ be the roots of the equation $x^2-x-1=0$. Define $T(n)=\frac{\alpha^n-\beta^n}{\sqrt5}$ $n\in N$. Which among the following holds TRUE?
\begin{qoptions}
\opt{A}{$T(2023)=T(2022)+T(2021)$}
\opt{B}{$T(10)=55$}
\opt{C}{$\sum_{n=1}^{2021}T(n)=T(2023)-1$}
\opt{D}{$\sum_{n=1}^{2021}T(n)=T(2023)+1$}
\end{qoptions}
\qn{7} $\displaystyle\lim_{x\to0}\left(\frac{1}{x^{2022}}\int_{x^{2023}}^{x^{2022}}\left(\frac{(1+t)^{\frac1{2021}}-1}{t}\right)dt\right)=\frac1P$\par
Then
\begin{qoptions}
\opt{A}{Number of positive divisors of $P$ is 4}
\opt{B}{$\ell n\,x=Px$ has exactly two real solutions}
\opt{C}{$2018^x+2019^x+2020^x=P^x$ has exactly one real root}
\opt{D}{Unit digit of $P^P$ is 1}
\end{qoptions}
\qn{8} Let $f:[0,1]\rightarrow(0,\infty)$ be an integrable function such that $f(1-x)=1/f(x)$ for all $x\in[0,1]$. The possible value of $\displaystyle\int_0^1f(x)\,dx$ can be
\begin{qoptions}
\opt{A}{$\frac{7}{22}$}
\opt{B}{$\pi$}
\opt{C}{$\frac1e$}
\opt{D}{$e$}
\end{qoptions}
\qn{9} An invertible function $f(x)=\sin^{-1}(x+2)+\cos^{-1}\left(\frac{x^2+x+1}{-x}\right)$ and $g$ is the inverse of function $f$, then\par
\textbf{Note :} Here $\mathrm{sgn}(x)$ denote signum function of $x$.
\begin{qoptions}
\opt{A}{$\mathrm{sgn}(f(g(x)))=1$}
\opt{B}{$\mathrm{sgn}(g(f(x)))=-1$}
\opt{C}{$g\left(\frac\pi2\right)=-1$}
\opt{D}{$f(x)$ is discontinuous at $x=-1$}
\end{qoptions}
\qn{10} Let $\begin{bmatrix}\frac12&\frac13&\frac16\\ \frac14&\frac15&\frac{11}{20}\\ \frac16&\frac17&\frac{29}{42}\end{bmatrix}=A$\par
If $A^K=\begin{bmatrix}a_K&b_K&c_K\\ d_K&e_K&f_K\\ g_K&h_K&i_K\end{bmatrix}$ then which of the following statements is/are true.
\begin{qoptions}
\opt{A}{$A^{2023}\begin{bmatrix}1\\1\\1\end{bmatrix}=\begin{bmatrix}\alpha\\\beta\\\gamma\end{bmatrix}$ where $\alpha,\beta,\gamma$ lies on the line of intersection of planes $3x+y+z=5$ and $2x-y+3z=4$}
\opt{B}{$\sum_{K=10}^{20}(a_K+b_K+c_K+d_K+e_K+f_K+g_K+h_K)=33$}
\opt{C}{$A^{10}\begin{bmatrix}x\\y\\z\end{bmatrix}=\begin{bmatrix}1\\1\\1\end{bmatrix}$ has no real solution}
\opt{D}{The number of $3\times3$ matrix $B$ with integer entries such that $AB-BA=\begin{bmatrix}1&1&1\\0&0&0\\0&1&0\end{bmatrix}$ is 2}
\end{qoptions}
\sect{SECTION-III}
\qn{1} Let $S$ be the sum of smallest element of each subset of the set $S=\left\{1,\frac12,\frac14,\ldots\ldots,\frac{1}{2^{2023}}\right\}$ (Ignore empty subset as it has no elements). Then the sum of digits of $S$ is \underline{\hspace{1cm}}.
\qn{2} Let $f(x)$ be a continuous function on $[0,1]$ such that $\displaystyle\int_0^1f(x)\,dx=\frac13+\int_0^1\left(f(x^2)\right)^2dx$. Then $10f\left(\frac14\right)$ equals \underline{\hspace{1cm}}.
\qn{3} If $I=\displaystyle\int_0^1\left(e^{x^2+x}+e^{x+\sqrt x}\right)dx$. Then $\ln(I+1)$ is \underline{\hspace{1cm}}.
\qn{4} Tangents are drawn from any point on straight line $y=x+8$ to the auxiliary circle of $x^2+16y^2=16$. If $A$ and $B$ are the points of contact of these tangent and $P$, $Q$ are corresponding points of $A$ and $B$ on the ellipse respectively. If locus of mid point of $PQ$ is $x^2+\alpha y^2+\beta x+\gamma y=0$ then value of $(\alpha-\beta+\gamma)$ is
\qn{5} Let $S$ be the solid tetrahedron with boundary points $(0,0,0)$, $(2,4,0)$, $(5,1,0)$, $(3,2,10)$.\par
Let $\alpha=\max\left\{q\ \middle|\ \left(\frac{12}{5},\frac{23}{10},q\right)\in S\right\}$\par
and $\beta=\max\left\{r\ \middle|\ \left(\frac92,\frac54,r\right)\in S\right\}$\par
Then $\frac{15}{\alpha-\beta}$ is \underline{\hspace{1cm}}.
\qn{6} Let $(x_0,y_0,z_0)$ be the solution of equation\par
$3x-4y+\frac{1}{xy}=12$\par
$4z-12x+\frac{1}{zx}=3$\par
$12y-3z+\frac{1}{yz}=4$\par
Then $\frac{z_0}{4x_0y_0}$ is \underline{\hspace{1cm}}.
\qn{7} Suppose $\hat u$ and $\hat v$ are two non-collinear unit vectors such that $\left|\frac{\hat u+\hat v}{2}+\hat u\times\hat v\right|=1$. The value of $\frac{|\hat u-\hat v|}{|\hat u\times\hat v|}$ is \underline{\hspace{1cm}}.
\qn{8} An unbiased coin is tossed $n$ times.\par
Define events $E_1$ and $E_2$ as follows :\par
$E_1$ : event that both heads and tails are present in $n$-tosses.\par
$E_2$ : event that coin shows up heads atmost once.\par
Find $n$ such that $E_1$ and $E_2$ are independent events.
\end{multicols}
\end{document}
